\documentclass[../../main.tex]{subfiles}
\begin{document}

\begin{problem}
\textbf{删数问题}
通过键盘输入一个高精度的正整数 n (n 的有效位数 ≤ 240),
去掉其中任意 s 个数字后, 剩下的数字按原左右次序将组成一个新的正整数.编程对给定的 n 和 s,
寻找一种方案, 使得剩下的数字组成的新数最小.
\end{problem}

\begin{solution}
    每一步总是选择一个使剩下的数最小的数字删去:
    即按高位到低位的顺序搜索，若各位数字递增，则删除最后一个数字，
    否则删除第一个递减区间的首字符，这样删一位便形成了一个新数字串。
    然后回到数字串首，按上述规则再删下一个数字。
    重复以上过程s次, 剩下的数字串便是问题的解.
    算法具体操作如下:


    \begin{algorithm}
        \caption{删除数字以获得最小数}
        \label{alg:getMinNumber}
        \begin{algorithmic}[1]
            \State \textbf{Input:} 字符串 \texttt{number}, 整数 \texttt{digitsToRemove}
            \State \textbf{Output:} 最小可能的数字
            \While{\texttt{digitsToRemove} > 0}
            \State 找到第一个比后一个大的数字位置 \texttt{pos}
            \State 删除 \texttt{number[pos]}
            \State \texttt{digitsToRemove} $\gets$ \texttt{digitsToRemove - 1}
            \EndWhile
            \State 删除前导零
            \If{\texttt{number}为空}
            \State \Return \texttt{"0"}
            \Else
            \State \Return \texttt{number}
            \EndIf
        \end{algorithmic}
    \end{algorithm}


    输入输出如下:
    \begin{figure}[htb]
        \centering
        \includegraphics[width=450pt]{figure/Figure_1.png}
        \caption{输入输出示例}
    \end{figure}
    \newpage
    该算法完整代码如下:
    \inputminted[
        frame=lines,
        framesep=2mm,
        baselinestretch=1.2,
        fontsize=\footnotesize
    ]{cpp}{project/project_3/src/main.cpp}
\end{solution}
\end{document}